\begin{proposition}[Newman 比值的代数证书与范数律（\texorpdfstring{$A_4$}{A4}）]\label{prop:pom-a4-newman-u4-algebraic-certificate}
\leanverified{paper\_pom\_a4\_newman\_u4\_algebraic\_certificate}
沿用注 \ref{rem:pom-collision-rh-break-a4} 的记号，令 \(r_4=\rho(A_4)\)，并令
\[
\Lambda_4:=\max_{\mu\in\mathrm{spec}(A_4)\setminus\{r_4\}}|\mu|.
\]
定义 Newman 比值与偏移量
\[
\upsilon_4:=\frac{\Lambda_4^2}{r_4},\qquad \delta_4:=\log \upsilon_4.
\]
则 \(\upsilon_4\) 为次数 \(20\) 的代数数，并且是多项式 \(Q_4(t)\) 在区间 \((1,1.01)\) 内的唯一实根；其中 \(Q_4(t)\) 的整系数证书为式 \eqref{eq:collision_kernel_A4_newman_u4_minpoly}。
此外其代数范数满足
\[
\mathrm{N}_{\QQ(\upsilon_4)/\QQ}(\upsilon_4)=16,
\]
从而 \(\delta_4\) 可被刚性化为“对数—代数根”的形式：\(\delta_4=\log(\upsilon_4)\)。
\par
并且 \(\upsilon_4\) 的极小多项式判别式的平方自由部分同为 \(-985219\)；
其整系数判别式分解证书见式 \eqref{eq:collision_kernel_A4_newman_u4_discriminant_factorization}。
\end{proposition}
\begin{proof}
由命题 \ref{prop:pom-a4-galois-s5}，特征多项式 \(\chi_{A_4}=\det(\lambda I-A_4)\) 在 \(\QQ\) 上不可约且其伽罗瓦群为 \(S_5\)。
令 \(\alpha_1,\dots,\alpha_5\) 为 \(\chi_{A_4}\) 的根，并取 \(\alpha_a=-\Lambda_4\) 为唯一负实根，\(\alpha_b=r_4\) 为 Perron 根。
则 \(\upsilon_4=\alpha_a^2/\alpha_b\) 属于分裂域；并且在 \(S_5\) 的作用下，其共轭集合恰为有序对轨道
\[
\mathcal{O}:=\left\{\frac{\alpha_i^2}{\alpha_j}:\ i\neq j\right\},
\]
其基数为 \(5\cdot 4=20\)，从而 \([\QQ(\upsilon_4):\QQ]=20\)。
对约束 \(\chi_{A_4}(\alpha_i)=\chi_{A_4}(\alpha_j)=0\) 与关系 \(t=\alpha_i^2/\alpha_j\) 作消元，可得到 \(t\) 的整系数最小多项式 \(Q_4(t)\)；其显式表达与标签见式 \eqref{eq:collision_kernel_A4_newman_u4_minpoly}。
\par
范数恒等式由首末系数给出：若 \(Q_4(t)=a_{20}t^{20}+\cdots+a_0\) 为 \(\upsilon_4\) 的最小多项式，则
\(
\mathrm{N}(\upsilon_4)=(-1)^{20}a_0/a_{20}=a_0/a_{20}=16
\)。
区间内唯一实根与数值一致性可由实根隔离算法（例如 Sturm 链）与高精度求根核验（脚本 \texttt{scripts/exp\_collision\_kernel\_A4\_newman\_u4\_certificate.py}）。证毕。
\end{proof}

% Auto-generated by scripts/exp_collision_kernel_A4_newman_u4_certificate.py
\begin{equation}\label{eq:collision_kernel_A4_newman_u4_minpoly}
\begin{aligned}
Q_4(t):={}&4t^{20} - 64t^{19} + 80t^{18} - 46260t^{17} - 127612t^{16} + 1584408t^{15} + 7293617t^{14} + 76445526t^{13} - 26190147t^{12} - 75382686t^{11} + 5220754t^{10} + 53887216t^{9} - 42120556t^{8} + 1739748t^{7} - 3577332t^{6} + 760256t^{5} + 144576t^{4} - 54720t^{3} - 10880t^{2} - 64t + 64.
\end{aligned}
\end{equation}

% Auto-generated by scripts/exp_collision_kernel_A4_newman_u4_discriminant_certificate.py
\begin{equation}\label{eq:collision_kernel_A4_newman_u4_discriminant_factorization}
\boxed{\;
\Disc\!\bigl(Q_4(t)\bigr)=-2^{110}\cdot 3^{6}\cdot 13^{2}\cdot 17^{8}\cdot 29^{6}\cdot 19309^{2}\cdot 985219^{7}\cdot 14623067^{2}\cdot 825912851^{2}\cdot 16389371777^{2}\cdot 80656775047^{2}\cdot 717880151227^{2}\cdot 1511624357228023^{2}\;
}
\end{equation}


\begin{remark}[轨道乘积的范数解释与 \texorpdfstring{$\det(A_4)$}{det(A4)} 的刚性联系]\label{rem:pom-a4-u4-orbit-product-norm}
沿用上证的轨道 \(\mathcal{O}\)，有精确恒等式
\[
\prod_{i\neq j}\frac{\alpha_i^2}{\alpha_j}
=\Bigl(\prod_{k=1}^5\alpha_k\Bigr)^{4}.
\]
又由 \(\prod_k\alpha_k=(-1)^5\chi_{A_4}(0)=-2\)（等价于 \(\det(A_4)=-2\)），故轨道乘积等于 \((-2)^4=16\)。
当 \(\Gal(\chi_{A_4}/\QQ)\cong S_5\) 时，\(\mathcal{O}\) 正是 \(\upsilon_4\) 的全体共轭，因而上式亦给出 \(\mathrm{N}_{\QQ(\upsilon_4)/\QQ}(\upsilon_4)=16\) 的“轨道乘积”解释。
\leanverified{prod\_numerator}
\leanverified{prod\_denominator}
\leanverified{orbit\_product\_identity}
\leanverified{orbit\_product\_n5}
\leanverified{paper\_pom\_a4\_u4\_orbit\_product\_norm}
\end{remark}
